\documentclass[aspectratio=169]{beamer}
\usetheme{Madrid}
\usecolortheme{default}

\usepackage{graphicx}
\usepackage{booktabs}
\usepackage{amsmath}

\title{How Accurate Are Learned World Models?}
\subtitle{An Empirical Analysis of DreamerV3's Imagination on Snake}
\author{Josh Manchester}
\institute{CS 5080 --- Reinforcement Learning\\Dr.\ Jugal Kalita\\Department of Computer Science\\University of Colorado Colorado Springs}
\date{February 19, 2026}

\begin{document}

% ---- SLIDE 1: Title ----
\begin{frame}
\titlepage
\end{frame}

% ---- SLIDE 2: Motivation ----
\begin{frame}{Model-Based RL: Learning by Dreaming}
\begin{itemize}
    \item Model-free RL learns by trial and error in the real environment
    \item Model-based RL learns \textbf{how the environment works} first, then plans
    \item DreamerV3 takes this further---it trains its policy \textbf{entirely inside the learned world model}
    \begin{itemize}
        \item The agent ``dreams'' about what might happen
        \item It never touches the real environment during policy optimization
    \end{itemize}
    \item One algorithm, fixed hyperparameters, masters Atari, robotics, and Minecraft
\end{itemize}

\vspace{0.5em}
\footnotesize{Hafner et al., \textit{Nature}, 2025}
\end{frame}

% ---- SLIDE 3: The Question ----
\begin{frame}{The Research Question}
\begin{center}
\Large
DreamerV3 trains inside imagined trajectories.\\[1em]
\textbf{But how accurate are those imaginations?}\\[1em]
\normalsize
If the world model predicts the wrong outcome for an action,\\
the agent is optimizing for a scenario that does not exist.\\[1em]
\textit{Does the agent need accurate dreams to play well?}
\end{center}
\end{frame}

% ---- SLIDE 4: Why This Matters ----
\begin{frame}{Why This Matters}
\begin{itemize}
    \item Most DreamerV3 research focuses on \textbf{end-task performance}---scores, win rates, learning speed
    \item Nobody has systematically measured \textbf{how faithful the imagination is} to reality
    \item Understanding imagination accuracy could reveal:
    \begin{itemize}
        \item Why model-based RL works despite imperfect models
        \item What the world model actually learns to represent
        \item Where model-based RL breaks down
    \end{itemize}
\end{itemize}
\end{frame}

% ---- SLIDE 5: Why Snake? ----
\begin{frame}{Why Snake?}
\begin{columns}
\begin{column}{0.55\textwidth}
\begin{itemize}
    \item \textbf{Simple, deterministic dynamics}
    \begin{itemize}
        \item Snake moves one cell per step
        \item Collisions end the game
        \item Only randomness: food placement
    \end{itemize}
    \item \textbf{Fully interpretable state}
    \begin{itemize}
        \item Snake head position
        \item Body segment locations
        \item Food location
    \end{itemize}
    \item \textbf{Easy to measure errors}---when the model imagines wrong, I can pinpoint exactly what went wrong
    \item Not yet studied with DreamerV3
\end{itemize}
\end{column}
\begin{column}{0.4\textwidth}
\centering
\textbf{Snake as an MDP}\\[0.5em]
\begin{tabular}{ll}
$\mathcal{S}$ & $N \times N$ grid \\
$\mathcal{A}$ & \{Up, Down, L, R\} \\
$R$ & +1 food, $-$1 collision \\
$\gamma$ & 0.999 \\
\end{tabular}
\end{column}
\end{columns}
\end{frame}

% ---- SLIDE 6: DreamerV3 Architecture ----
\begin{frame}{DreamerV3 Architecture Overview}
\begin{itemize}
    \item \textbf{World Model} (RSSM---Recurrent State-Space Model):
    \begin{itemize}
        \item Encodes observations into latent states
        \item Deterministic path (GRU) captures persistent information
        \item Stochastic path models environment uncertainty
        \item Predicts next states, rewards, and reconstructed observations
    \end{itemize}
    \item \textbf{Actor-Critic} trained entirely in imagined trajectories:
    \begin{itemize}
        \item Actor selects actions in latent space
        \item Critic estimates value of imagined states
        \item No real environment interaction during policy learning
    \end{itemize}
    \item Key innovations: symlog predictions, fixed hyperparameters across domains
\end{itemize}

\vspace{0.5em}
\footnotesize{Built on: Hafner et al. (2019, 2020, 2021, 2025)}
\end{frame}

% ---- SLIDE 7: Method - Training & Extraction ----
\begin{frame}{Method: Training and Trajectory Extraction}
\textbf{Step 1: Train DreamerV3 on Snake}
\begin{itemize}
    \item Gymnasium-compatible Snake environment, 64$\times$64 RGB observations
    \item 1 million environment steps, checkpoints every 50K steps (20 snapshots)
\end{itemize}

\vspace{0.5em}
\textbf{Step 2: Extract Matched Trajectories at Each Checkpoint}
\begin{enumerate}
    \item Record a real trajectory using the current policy
    \item Feed the initial observation into the world model
    \item Roll out the world model with the \textbf{same action sequence}
    \item Compare real vs.\ imagined frames side by side
\end{enumerate}

\vspace{0.5em}
Same actions $\rightarrow$ only difference is the world model's prediction accuracy.
\end{frame}

% ---- SLIDE 8: Metrics ----
\begin{frame}{Measuring Imagination Accuracy}
\begin{columns}
\begin{column}{0.48\textwidth}
\textbf{Frame-Level Metrics}
\begin{itemize}
    \item \textbf{MSE}---pixel-level error
    \item \textbf{SSIM}---structural fidelity
\end{itemize}
\vspace{0.5em}
Captures overall visual quality of the world model's predictions.
\end{column}
\begin{column}{0.48\textwidth}
\textbf{Semantic Metrics}
\begin{itemize}
    \item Head position error
    \item Body segment accuracy
    \item Food location accuracy
    \item Collision prediction
\end{itemize}
\vspace{0.5em}
Captures whether the model understands the \textit{game state}, not just the pixels.
\end{column}
\end{columns}

\vspace{1em}
\begin{center}
\textbf{Key insight:} Blurry but correct $\neq$ sharp but wrong.\\
Separating these tells me what the model actually learns.
\end{center}
\end{frame}

% ---- SLIDE 9: Correlation Analysis ----
\begin{frame}{Correlation: Do Better Dreams Mean Better Play?}
\begin{itemize}
    \item Track both \textbf{imagination accuracy} and \textbf{game performance} across 20 checkpoints
    \item Three possible outcomes:
\end{itemize}

\vspace{0.5em}
\begin{tabular}{lp{9cm}}
\toprule
\textbf{Outcome} & \textbf{What It Means} \\
\midrule
Accuracy $\uparrow$, Performance $\uparrow$ & Better dreams drive better play---imagination accuracy is critical \\
Accuracy flat, Performance $\uparrow$ & Agent succeeds despite imperfect dreams---policy is robust to model errors \\
Accuracy $\uparrow$, Performance flat & Model improves but policy does not benefit---something else is the bottleneck \\
\bottomrule
\end{tabular}

\vspace{0.5em}
I will also test which \textbf{semantic metrics} (food location, collision prediction) are most predictive of performance.
\end{frame}

% ---- SLIDE 10: Timeline ----
\begin{frame}{Timeline}
\begin{table}
\centering
\begin{tabular}{ll}
\toprule
\textbf{Date} & \textbf{Milestone} \\
\midrule
Feb 19 & Proposal presentation and paper \\
Mar 1 & Snake environment, DreamerV3 training \\
Mar 15 & Imagination extraction pipeline \\
Mar 31 & Checkpoint analysis complete \\
Apr 1 & Midterm presentation and paper \\
Apr 15 & Correlation and semantic metrics \\
May 7 & Final presentation \\
May 12 & Final paper \\
\bottomrule
\end{tabular}
\end{table}

\vspace{0.5em}
\textbf{Hardware:} RTX 5070 Ti (16GB VRAM), PyTorch 2.7.0, CUDA 12.8
\end{frame}

% ---- SLIDE 11: Questions ----
\begin{frame}
\begin{center}
\Large \textbf{Questions?}\\[2em]
\normalsize
Josh Manchester\\
\texttt{jmanches@uccs.edu}
\end{center}
\end{frame}

\end{document}
